% Sección 1
\section{Para qué este informe y cómo leerlo}

Este informe responde a un encargo de la mentoría de la hackathon. Los mentores nos pidieron conocer el problema a fondo antes de defender una solución, y dejaron cinco preguntas. Quieren saber por qué hacemos esto y para qué sirve, y a quién le sirve y a quién impacta. También preguntaron si existe algo parecido y por qué no se usa en Manizales o en la región. Las dos últimas miran hacia adelante y piden el prototipo que medirá el efecto en la región y el piloto de 30 a 90 días. La expresión que pusieron en el centro fue «recuperación del territorio». Con ella señalaron que la respuesta tiene que hablar de personas y de lugares además de tecnología.

El documento está escrito para quien no vive en la región. Reúne lo que el equipo ya investigó y lo ordena alrededor de esas preguntas. Es un documento vivo. Esta versión 0.1 recoge los informes de investigación que hay en el repositorio, y las siguientes añadirán lo que traigan otros agentes durante la noche.

\subsection{Niveles de evidencia}

Cada cifra y cada afirmación del informe lleva uno de cuatro niveles de evidencia. Las tablas los abrevian con la palabra en negrilla.

\begin{description}[leftmargin=1.4em, style=nextline, itemsep=0.35em]
\item[\textsf{Oficial}.] Fuente oficial. Es un documento de la entidad que produce el dato o un decreto. Por ejemplo, el Decreto 1261 de 2026 cita el balance de la Unidad Nacional para la Gestión del Riesgo (UNGRD).
\item[\textsf{Prensa}.] Prensa que cita a una entidad. Es una nota de prensa, o el documento de un tercero, que reproduce lo que dijo una entidad o un funcionario. Una cifra que la investigación del equipo calificó con confianza media queda en este nivel aunque venga de un documento institucional, como los informes de Asocapitales. Así no subimos el nivel que le dio la investigación.
\item[\textsf{Inferencia}.] Inferencia nuestra. Es una conclusión del equipo a partir de los datos, o un resultado de nuestra simulación. Las cifras de simulación salen de un modelo con supuestos propios y no de una operación real.
\item[\textsf{Hipótesis}.] Hipótesis por validar. Es algo que creemos y que nadie ha confirmado todavía. Se valida hablando con quien estuvo en la operación.
\end{description}

Las tablas añaden entre paréntesis la confianza que asignó el informe de investigación del equipo, sin cambios. «Alta» corresponde a una entidad oficial o un decreto. «Media» corresponde a prensa con fuente citada o a una cifra preliminar. Cuando el dato viene de un informe que no calificó la confianza, la tabla dice «sin calificar».

\subsection{La regla para llevar una cifra al pitch}

El equipo tiene una regla. Ninguna cifra va al pitch hasta que una persona abra la fuente y la anote en \ruta{docs/FUENTES.md}. Hoy ninguna cifra de este informe cumple esa condición. Los datos vienen de informes que prepararon agentes de investigación el 20 de noviembre de 2026, y nadie del equipo ha abierto todavía los enlaces. Varias consultas web de esos agentes quedaron bloqueadas, de modo que hay vacíos que este informe señala sin rellenarlos.

Las cifras del sismo son preliminares y cambian con cada corte. Por eso cada una aparece con su fecha, y cuando dos fuentes difieren se dan las dos. El informe usa además tres marcas. \pendiente{} indica una respuesta que nadie del equipo ha conseguido todavía. \eninvestigacion{} indica algo que dos agentes están buscando esta noche. \porverificar{} señala un dato que nos pareció dudoso y que dejamos tal como está en la fuente.

Las fuentes de cada dato están en los informes de investigación del repositorio, en la carpeta \ruta{docs/investigacion/}. Los que más usa este informe son el de evidencia del problema (\ruta{04_evidencia_del_problema.md}) y el de referentes latinoamericanos (\ruta{02_referentes_latinoamerica.md}). El anexo reúne en una tabla las fuentes que usa este informe.
